% =============================================================================
%  CHAPITRE 2 : MÉTHODOLOGIE DE L'ÉTUDE
% =============================================================================
\entete{Chapitre 2 : Méthodologie de l'étude}
\chapter{Méthodologie de l'étude}
\label{chap:methodo}

\introchap

La valeur d'un résultat tient à la rigueur du chemin qui y mène. Ce chapitre
expose donc la démarche suivie, avec le détail nécessaire à sa reproduction. Il
précise le type de recherche entreprise, la manière dont le phénomène étudié a
été rendu mesurable, les personnes auprès desquelles les données ont été
recueillies, les outils employés pour les recueillir et les analyser, puis la
méthode de conduite du projet informatique. Il se clôt sur les difficultés
rencontrées et les limites assumées, qui bornent la portée des résultats.

% =============================================================================
\section{Nature de l'étude et approche méthodologique}
\label{sec:nature}

\subsection{Niveau et type de recherche}

Le choix du niveau de recherche détermine la profondeur du traitement
scientifique et, par voie de conséquence, les méthodes admissibles. Quatre
niveaux sont habituellement distingués : perceptuel, appréhensif, compréhensif
et intégrateur.

La présente étude relève du \textbf{niveau intégrateur}. Nous ne nous sommes en
effet pas limité à décrire une situation ni à en expliquer les causes : nous
sommes intervenu sur la réalité observée en concevant, en construisant et en
mettant en service un dispositif destiné à la transformer. Le type d'étude
correspondant est la \textbf{recherche interactive}, également désignée dans la
littérature en systèmes d'information sous le nom de \emph{recherche par la
conception} (Hevner et al., 2004). Elle se caractérise par :

\begin{itemize}
  \item une forte interaction avec les acteurs du terrain : le personnel de la
        Fondation a été associé à chaque étape ;
  \item une démarche participative, les besoins ayant été formulés, corrigés et
        arbitrés avec l'encadreur professionnel ;
  \item la résolution d'un problème réel et non d'un cas d'école ;
  \item le recours au prototypage : chaque fonctionnalité a été construite,
        soumise, éprouvée, puis corrigée.
\end{itemize}

Ce niveau excède ce qui est exigé d'un travail de licence, lequel peut se
limiter au niveau perceptuel ou appréhensif. Ce choix se justifie par la nature
même de la commande : la Fondation n'attendait pas un diagnostic, elle attendait
un outil.

\subsection{Approche méthodologique}

Nous avons retenu une \textbf{approche mixte}, associant les apports qualitatifs
et quantitatifs, pour trois raisons.

\begin{enumerate}
  \item La compréhension du besoin, ce qu'une institution patrimoniale attend
        d'une visite à distance, appelle des données qualitatives : entretiens,
        observation, analyse des documents internes.
  \item L'évaluation du dispositif construit appelle au contraire des données
        quantitatives : temps de réponse, taux de réussite, scores de
        correspondance, coûts mesurés.
  \item Le croisement des deux permet de vérifier qu'un dispositif
        techniquement performant est également jugé utile et utilisable par
        ceux à qui il est destiné.
\end{enumerate}

% =============================================================================
\section{Variables et indicateurs de l'étude}
\label{sec:variables}

\subsection{Identification des variables}

Conformément à la méthode, les variables de l'étude sont extraites de
l'hypothèse : les composantes énoncées dans l'hypothèse générale en constituent
les variables explicatives, et l'effet attendu en constitue la variable
expliquée. Chaque variable explicative correspond ainsi à une hypothèse
spécifique, et une seule.

La \textbf{variable expliquée} est l'effet recherché :

\begin{itemize}
  \item \textbf{la qualité de l'expérience de visite à distance}, entendue comme
        la capacité du dispositif à procurer au visiteur éloigné une perception
        des lieux, une compréhension des œuvres et une possibilité
        d'interrogation comparables à celles d'une visite sur place.
\end{itemize}

Les quatre \textbf{variables explicatives} sont les composantes qui produisent
cet effet :

\begin{itemize}
  \item \textbf{Variable explicative 1 : la restitution spatiale des espaces et
        des œuvres}, c'est-à-dire les panoramas à 360 degrés reliés entre eux,
        la présentation en trois dimensions et la réalité augmentée à taille
        réelle (hypothèse spécifique 1) ;
  \item \textbf{Variable explicative 2 : le rapprochement documentaire avec les
        collections ouvertes}, c'est-à-dire la recherche automatique des œuvres
        apparentées et leur validation par le conservateur (hypothèse spécifique
        2) ;
  \item \textbf{Variable explicative 3 : l'ancrage documentaire de l'assistant
        conversationnel}, c'est-à-dire la construction des réponses à partir des
        seuls documents validés par la Fondation (hypothèse spécifique 3) ;
  \item \textbf{Variable explicative 4 : l'isolation des données au niveau de la
        base}, c'est-à-dire le mécanisme qui empêche une organisation de
        consulter les données d'une autre, et dont dépend la possibilité
        d'accueillir de nouvelles institutions sans modifier le code (hypothèse
        spécifique 4).
\end{itemize}

\subsection{Opérationnalisation : les indicateurs}

Une variable n'est exploitable que si l'on sait la mesurer. Le
tableau~\ref{tab:variables} présente, pour chaque variable, les indicateurs
retenus, leur unité, la source de la mesure et, dans une dernière colonne, le
statut de l'indicateur au terme du stage : \emph{mesuré}, lorsqu'une valeur a
effectivement été relevée et figure au chapitre~3, ou \emph{non mesuré} dans le
cas contraire. Cette dernière colonne n'est pas une précaution de style : sans
elle, un indicateur annoncé se confondrait avec un indicateur renseigné, et le
lecteur ne pourrait pas apprécier ce que le travail établit réellement.

\begin{table}[htbp]
\centering
\caption{Variables, indicateurs, sources de mesure et statut}
\label{tab:variables}
\scriptsize
\begin{tabularx}{\textwidth}{|L{2.2cm}|X|L{1.7cm}|L{2.3cm}|C{1.8cm}|}
\hline
\rowcolor{keycetete}
\thead{Variable} & \thead{Indicateurs} & \thead{Unité} & \thead{Source de la mesure} & \thead{Statut} \\ \hline
\multirow{5}{*}{\parbox{2.2cm}{\textbf{Variable 1}\\ Restitution spatiale}}
 & Nombre de salles restituées en panorama & effectif & Base de données & Non mesuré \\ \cline{2-5}
& Nombre de points d'intérêt cliquables & effectif & Base de données & Non mesuré \\ \cline{2-5}
& Délai d'affichage d'une salle & secondes & Mesure navigateur & Non mesuré \\ \cline{2-5}
& Écart dimensionnel du modèle 3D & centimètres & Mesure sur maillage & \textbf{Mesuré} \\ \cline{2-5}
& Couverture des appareils testés & \% & Essais sur appareils & \textbf{Mesuré} \\ \hline
\multirow{5}{*}{\parbox{2.2cm}{\textbf{Variable 2}\\ Mise en relation documentaire}}
 & Nombre de collections interrogées & effectif & Journal des appels & \textbf{Mesuré} \\ \cline{2-5}
& Nombre d'œuvres apparentées identifiées & effectif & Table des rapprochements & \textbf{Mesuré} \\ \cline{2-5}
& Score de correspondance & /100 & Algorithme de notation & \textbf{Mesuré} \\ \cline{2-5}
& Taux de validation par le conservateur & \% & Table des rapprochements & Non mesuré \\ \cline{2-5}
& Nombre de pays de dispersion & effectif & Données de dispersion & \textbf{Mesuré} \\ \hline
\multirow{5}{*}{\parbox{2.2cm}{\textbf{Variable 3}\\ Ancrage documentaire de l'assistant}}
 & Taux de réponses appuyées sur un document interne & \% & Journal des questions & \textbf{Mesuré} \\ \cline{2-5}
& Taux d'aveu d'ignorance en l'absence de source & \% & Grille d'évaluation & Non mesuré \\ \cline{2-5}
& Taux de refus des questions hors périmètre & \% & Grille d'évaluation & Non mesuré \\ \cline{2-5}
& Délai moyen de réponse & secondes & Journal serveur & Non mesuré \\ \cline{2-5}
& Questions restées sans réponse & effectif & Journal des questions & \textbf{Mesuré} \\ \hline
\multirow{4}{*}{\parbox{2.2cm}{\textbf{Variable 4}\\ Cloisonnement et ouverture}}
 & Tables protégées par une règle de sécurité & effectif & Schéma de la base & \textbf{Mesuré} \\ \cline{2-5}
& Tentatives d'accès inter-organisations refusées & \% & Essais d'intrusion & Non mesuré \\ \cline{2-5}
& Délai de création d'une organisation & minutes & Essai de bout en bout & Non mesuré \\ \cline{2-5}
& Lignes de code à modifier pour un nouvel arrivant & effectif & Analyse du code & \textbf{Mesuré} \\ \hline
\end{tabularx}
\source{Élaboré par nos soins.}
\end{table}

La colonne \emph{Statut} appelle une lecture franche : sur les dix-neuf
indicateurs retenus, \textbf{onze ont effectivement produit une valeur} et huit
n'en ont pas produit. Les indicateurs mesurés sont ceux qui portent les quatre
hypothèses de l'étude, et leurs valeurs sont rassemblées au
tableau~\ref{tab:mesures}. Les huit autres relèvent de deux situations
distinctes. Les uns supposaient des médias qui n'ont pas pu être produits durant
le stage, panoramas de salles réelles et points d'intérêt associés, pour les
raisons exposées à la section~\ref{sec:difficultes}. Les autres, délais
d'affichage et de réponse, taux de validation, essais d'intrusion, auraient exigé
un protocole de mesure répété que la durée du stage n'a pas permis d'instruire.
Les annoncer sans les avoir mesurés aurait été plus commode ; les déclarer non
mesurés délimite exactement la portée des conclusions que ce travail autorise.

% =============================================================================
\section{Population et échantillon}
\label{sec:echantillon}

\subsection{Population de l'étude}

La population de l'étude est constituée de trois ensembles d'acteurs concernés
par la valorisation numérique du patrimoine de la Fondation Jean Félicien Gacha
au cours de la période du 20~juin au 20~août 2026 :

\begin{itemize}
  \item \textbf{le personnel de la Fondation et du \emph{Think Tank} Fou'ndjem},
        c'est-à-dire les personnes intervenant dans la gestion des espaces, des
        collections, de l'accueil et de la communication ;
  \item \textbf{les visiteurs de la Fondation}, présents sur le site durant la
        période, âgés de dix-huit ans et plus.
\end{itemize}

\subsection{Technique d'échantillonnage et taille de l'échantillon}

Compte tenu de la taille réduite et bien identifiée des populations concernées,
un échantillonnage probabiliste n'aurait apporté aucun gain de représentativité.
Nous avons donc retenu deux techniques \textbf{non probabilistes} :

\begin{itemize}
  \item \textbf{le choix raisonné} pour le personnel : ont été retenues les
        personnes exerçant effectivement une responsabilité sur les espaces, les
        collections ou la communication, seules à même de renseigner les
        pratiques réelles ;
  \item \textbf{la convenance} pour les visiteurs : ont été sollicitées les
        personnes présentes sur le site durant les journées d'observation et
        ayant accepté de répondre.
\end{itemize}

Le tableau~\ref{tab:echantillon} présente la composition de l'échantillon.

\begin{table}[htbp]
\centering
\caption{Composition de l'échantillon de l'étude}
\label{tab:echantillon}
\small
\begin{tabularx}{\textwidth}{|X|C{2.4cm}|C{2.6cm}|L{3.4cm}|}
\hline
\rowcolor{keycetete}
\thead{Catégorie d'acteurs} & \thead{Effectif} & \thead{Technique} & \thead{Outil de collecte} \\ \hline
Personnel de la Fondation et du \emph{Think Tank} Fou'ndjem & 06 & Choix raisonné & Entretien semi-directif \\ \hline
Visiteurs présents sur le site & 42 & Convenance & Questionnaire \\ \hline
\rowcolor{keycetotal}
\textbf{Total} & \textbf{48} & — & — \\ \hline
\end{tabularx}
\source{Enquête de terrain, juin--août 2026.}
\end{table}

Cet échantillon ne prétend pas à la représentativité statistique. Il vise la
\emph{saturation de l'information} : les entretiens ont été conduits jusqu'à ce
que les échanges cessent d'apporter des éléments nouveaux, critère usuel en
recherche qualitative.

% =============================================================================
\section{Outils de collecte et techniques d'analyse des données}
\label{sec:outils}

\subsection{Outils de collecte}

Quatre outils ont été mobilisés, chacun rattaché à un type de donnée précis.

\begin{enumerate}
  \item \textbf{Le guide d'entretien semi-directif} (annexe~\ref{ann:guide}),
        organisé autour de quatre thèmes : pratiques de valorisation,
        difficultés, attentes et contraintes de moyens. Sa forme semi-directive
        s'est révélée décisive : la question de la dispersion des œuvres
        apparentées n'était pas prévue et a émergé d'un entretien.

  \item \textbf{Le questionnaire} (annexe~\ref{ann:questionnaire}), adressé aux
        visiteurs : dix-huit questions fermées et deux ouvertes, en quatre
        rubriques. Un pré-test auprès de cinq personnes a conduit à reformuler
        trois questions ambiguës et à en supprimer une.

  \item \textbf{La grille d'observation}, utilisée lors des visites guidées pour
        relever le parcours effectif des visiteurs, leurs points d'arrêt et les
        questions spontanément posées au guide. Ces relevés ont orienté la
        conception du parcours immersif.

  \item \textbf{La fiche documentaire}, qui a encadré le dépouillement des
        inventaires, notices et plaquettes de la Fondation, ainsi que la revue
        de la littérature scientifique.
\end{enumerate}

\subsection{Techniques d'analyse}

Les entretiens ont fait l'objet d'une \textbf{analyse thématique de contenu} :
après retranscription, les propos ont été découpés en unités de sens, codés,
puis regroupés en catégories, d'où ont émergé cinq thèmes, l'accessibilité, la
perception de l'espace, la mise en relation des objets, la fiabilité de
l'information et la pérennité économique du dispositif. Aucun logiciel
spécialisé n'a été employé, le volume traité ne le justifiant pas et un
dépouillement manuel offrant une maîtrise supérieure du codage. Les réponses au
questionnaire ont pour leur part été traitées par \textbf{statistique
descriptive}, sous forme d'effectifs, de fréquences, de pourcentages et de
moyennes présentés en tableaux et en graphiques ; les mesures relevées sur le
système, délais, scores et taux, ont fait l'objet de comparaisons avant et après
traitement, notamment pour évaluer l'apport de l'enrichissement des requêtes
documentaires. Le tableur a suffi à ces calculs.

L'assistant conversationnel a été évalué séparément, au moyen d'une
\textbf{grille d'évaluation} appliquée à un jeu de questions construit à
l'avance (annexe~\ref{ann:grille}). Ce jeu comporte quatre familles : des
questions \textbf{documentées}, dont la réponse figure dans les textes fournis
par la Fondation et auxquelles l'assistant doit répondre exactement ; des
questions \textbf{non documentées}, portant sur le périmètre de la Fondation
mais sans source disponible, face auxquelles il doit reconnaître qu'il ne sait
pas ; des questions \textbf{hors périmètre}, comme la météo ou l'actualité,
qu'il doit refuser poliment ; et des questions \textbf{pièges}, formulées de
manière à suggérer une réponse fausse, qu'il ne doit pas suivre. Cette
construction est délibérée : elle mesure non seulement ce que l'assistant sait,
mais surtout ce qu'il refuse d'inventer, critère décisif en médiation
culturelle.
\section{Méthode de conduite du projet}
\label{sec:devops}

La conduite du projet a reposé sur deux méthodes complémentaires, qui ne
répondent pas à la même question. La première, le \textbf{Scrumban}, organise le
travail : elle dit dans quel ordre les tâches sont traitées et à quel rythme
elles sont présentées à l'encadreur. La seconde, le \textbf{DevOps}, organise la
livraison : elle dit comment une modification écrite parvient jusqu'aux
utilisateurs sans casser ce qui fonctionnait.

\subsection{Le Scrumban comme méthode de travail}

Les méthodes agiles ont remplacé la planification détaillée par des cycles
courts au terme desquels une version utilisable est présentée. Deux d'entre
elles dominent. Le \textbf{Scrum} découpe le travail en sprints de durée fixe,
avec des rôles et des réunions ritualisées. Le \textbf{Kanban} supprime les
sprints au profit d'un flux continu : les tâches circulent sur un tableau, de
« à faire » à « terminé », et une limite est posée sur le nombre de tâches
menées de front, afin d'éviter que tout soit commencé et rien achevé.

Le \textbf{Scrumban} combine les deux : le flux continu et la limitation du
travail en cours viennent du Kanban, la cadence régulière de revue vient du
Scrum. C'est cette méthode que nous avons retenue, pour quatre raisons propres
au contexte du stage.

\begin{enumerate}
  \item \textbf{Le projet était conduit par une seule personne.} Les rituels du
        Scrum, mêlée quotidienne, estimation collective, revue de sprint en
        équipe, n'ont pas de sens à un seul développeur : ils consommeraient du
        temps sans produire de coordination.

  \item \textbf{Le besoin a évolué en cours de route.} La question du
        rapprochement documentaire n'était pas prévue et a émergé d'un
        entretien. Un sprint figé aurait obligé à attendre la fin du cycle pour
        l'intégrer ; le flux continu du Kanban permet d'insérer une tâche dès
        qu'elle est priorisée.

  \item \textbf{Une cadence de revue restait nécessaire.} L'encadreur
        professionnel suivait un programme hebdomadaire. Conserver du Scrum une
        revue chaque semaine, alignée sur ce programme, donnait un rendez-vous
        régulier de démonstration et d'arbitrage. Le chronogramme présenté au
        chapitre~3 en est la trace directe.

  \item \textbf{Le risque principal était la dispersion.} Le projet comporte
        vingt fonctionnalités hétérogènes. La limite de travail en cours du
        Kanban, fixée à deux tâches, a été le garde-fou le plus utile : elle
        oblige à terminer avant d'entreprendre.
\end{enumerate}

\begin{figure}[htbp]
\centering
\schema{scrumban}
\caption{Le tableau Scrumban adopté pour la conduite du projet}
\label{fig:scrumban}
\source{Élaboré par nos soins.}
\end{figure}

\subsection{Le DevOps comme démarche de livraison}

La construction d'un logiciel ne s'achève pas à l'écriture du code : encore
faut-il le livrer, le corriger et le faire évoluer sans interrompre le service.
Trois raisons ont conduit à adopter une démarche DevOps dès le début du projet
plutôt qu'à la fin.

\begin{enumerate}
  \item \textbf{La sécurité du travail.} Un stage de deux mois ne laisse aucune
        marge pour reconstruire ce qui aurait été cassé sans qu'on s'en
        aperçoive. Des contrôles automatiques déclenchés à chaque modification
        signalent la régression en quelques minutes.

  \item \textbf{La reproductibilité.} La Fondation devra un jour reprendre le
        système. Une mise en production qui repose sur une suite de gestes
        manuels mémorisés par une seule personne n'est pas transmissible ; une
        mise en production décrite dans un fichier l'est.

  \item \textbf{La sensibilité du dispositif aux détails d'exploitation.}
        L'expérience du projet a montré qu'un modèle tridimensionnel servi avec
        une étiquette de type erronée s'affiche correctement dans le navigateur
        mais empêche toute session de réalité augmentée. Ce type d'anomalie ne
        se voit pas à l'œil : il se vérifie automatiquement, ou il passe
        inaperçu jusqu'au jour de la démonstration.
\end{enumerate}

\subsection{Les pratiques retenues}

Quatre pratiques ont été mises en œuvre :

\begin{itemize}
  \item \textbf{le contrôle de version} de l'intégralité du projet, code
        applicatif et description de l'infrastructure comprise ;
  \item \textbf{l'intégration continue} : chaque dépôt de modification déclenche
        la compilation et une batterie de contrôles portant sur le résultat
        produit ;
  \item \textbf{le déploiement continu} : une modification acceptée est publiée
        automatiquement, avec purge du cache de diffusion et vérification, après
        coup, de ce que reçoit réellement le visiteur ;
  \item \textbf{l'infrastructure décrite par le code}, afin que l'environnement
        de production puisse être recréé à l'identique.
\end{itemize}

La chaîne effectivement mise en place et son évolution proposée sont décrites
respectivement au chapitre~3 et au chapitre~4.

% =============================================================================
\section{Outils et environnement de l'étude}
\label{sec:outilsetude}

Le tableau~\ref{tab:outils} recense les outils employés, en distinguant leur
rôle dans le travail.

\begin{table}[htbp]
\centering
\caption{Outils, technologies et environnements mobilisés}
\label{tab:outils}
\scriptsize
\begin{tabularx}{\textwidth}{|L{2.5cm}|L{4.3cm}|X|}
\hline
\rowcolor{keycetete}
\thead{Domaine} & \thead{Outil, technologie et version} & \thead{Rôle dans le travail} \\ \hline
\multirow{4}{*}{\parbox{2.5cm}{Interface}}
 & Vue 3.5 & Cadriciel de construction de l'interface \\ \cline{2-3}
& Vite 5.4 & Chaîne de compilation et serveur de développement \\ \cline{2-3}
& PrimeVue 4.2 & Bibliothèque de composants d'interface pour l'ERP \\ \cline{2-3}
& Pinia 2.3, Vue Router 4.5, Vue I18n 9.14 & État applicatif, navigation, bilinguisme \\ \hline
\multirow{4}{*}{\parbox{2.5cm}{Données et services}}
 & PostgreSQL 17.6 & Base de données relationnelle \\ \cline{2-3}
& Supabase (client JS 2.110) & Authentification, sécurité au niveau de la ligne, fonctions serveur \\ \cline{2-3}
& \texttt{pgvector} (index HNSW) & Stockage et recherche des plongements \\ \cline{2-3}
& Deno 2 & Exécution des fonctions serveur \\ \hline
\multirow{3}{*}{\parbox{2.5cm}{Immersion et 3D}}
 & WebGL 2 via three.js 0.183 (moteur écrit sur mesure) & Rendu des panoramas à 360~degrés et du relief \\ \cline{2-3}
& \texttt{model-viewer} 4.0, WebXR & Affichage 3D et réalité augmentée \\ \cline{2-3}
& Meshroom 2023.3, Blender 4.2 & Photogrammétrie et préparation des maillages \\ \hline
\multirow{4}{*}{\parbox{2.5cm}{Intelligence artificielle}}
 & \texttt{gemini-3.6-flash} (Google) & Modèle principal de génération et de reformulation, retenu dans sa déclinaison \emph{flash}, la plus légère de la famille \\ \cline{2-3}
& \texttt{llama3.2:1b} (Ollama, exécution locale) & Modèle de base employé en développement, hors service distant \\ \cline{2-3}
& \texttt{gte-small} (384 dimensions) & Production des plongements lexicaux \\ \cline{2-3}
& ElevenLabs (API v1) & Synthèse vocale du guide \\ \hline
\multirow{6}{*}{\parbox{2.5cm}{Industrialisation}}
 & Git 2.4x, GitHub & Contrôle de version et hébergement du dépôt \\ \cline{2-3}
& GitHub Actions & Intégration et déploiement continus \\ \cline{2-3}
& Docker 27 & Conteneurisation de l'application \\ \cline{2-3}
& Terraform $\geq$ 1.6, fournisseur AWS 5.60 & Description de l'infrastructure \\ \cline{2-3}
& Infracost 0.10.42 & Estimation du coût de l'infrastructure à partir de sa description \\ \cline{2-3}
& AWS (S3, CloudFront, Route~53, ACM, IAM) & Hébergement et diffusion \\ \hline
\multirow{3}{*}{\parbox{2.5cm}{Conception et rédaction}}
 & Visual Studio Code & Édition du code \\ \cline{2-3}
& \texttt{draw.io}, Ti\emph{k}Z & Réalisation des schémas \\ \cline{2-3}
& \LaTeX{} (MiKTeX, moteur pdf\TeX) & Composition du présent rapport \\ \hline
\end{tabularx}
\source{Élaboré par nos soins ; versions relevées dans les fichiers de
dépendances du projet au 4 septembre 2026.}
\end{table}

Deux précisions s'imposent sur les modèles de langue. La première tient à leur
instabilité : les identifiants de modèles sont périssables, plusieurs ayant été
retirés des catalogues des fournisseurs au cours même du stage. Le système ne
les inscrit donc jamais dans le code : ils sont lus dans une variable
d'environnement, ce qui permet d'en changer sans redéployer l'application. La
seconde tient au choix de la déclinaison : \texttt{gemini-3.6-flash} est la
variante la plus légère de sa famille, retenue précisément pour cela, le budget
de l'institution ne permettant pas d'appeler un modèle de raisonnement à chaque
question de visiteur.

% =============================================================================
\section{Difficultés rencontrées et limites méthodologiques}
\label{sec:difficultes}

L'honnêteté méthodologique commande d'exposer les obstacles rencontrés, car ils
conditionnent la portée des résultats.

\begin{enumerate}
  \item \textbf{Un format de modélisation 3D imposé par Apple.} L'application
        utilisée pour modéliser les objets en trois dimensions ne produisait
        qu'un seul format de sortie, l'USDZ, propre à l'écosystème d'Apple. Or
        les navigateurs et les appareils Android ne le lisent pas : en l'état,
        la moitié du public visé n'aurait rien vu. La difficulté a été levée en
        introduisant une étape de conversion vers le format GLB, qui est le
        format d'échange du web. Chaque œuvre dispose donc désormais de ses deux
        versions, l'USDZ pour les appareils Apple et le GLB pour les autres.

  \item \textbf{La mise en place des sous-domaines par institution.} Pour qu'une
        institution qui rejoint la plateforme dispose de sa propre adresse, il
        fallait faire exister un sous-domaine par organisation. Le montage s'est
        révélé plus délicat que prévu : il exige un certificat de sécurité
        couvrant à la fois le domaine principal et l'ensemble de ses
        sous-domaines, obtenu auprès du service de gestion des certificats
        d'AWS, ainsi que des enregistrements à caractère générique dans la zone
        hébergée de Route 53. La validation du certificat passe par le système
        de noms de domaine public : tant que la délégation du domaine n'est pas
        effective, la demande reste en attente puis expire, sans message
        explicite. C'est la difficulté qui a demandé le plus de tentatives.

  \item \textbf{Absence de clés d'accès à quatre grandes collections.} Europeana,
        le Rijksmuseum, la Smithsonian Institution et les Harvard Art Museums
        exigent une clé que nous n'avons pu obtenir. Le rapprochement porte donc
        sur cinq sources et non sur neuf, ce qui minore le nombre d'œuvres
        apparentées trouvées.

  \item \textbf{Comparaison portant sur le texte et non sur l'image.} Faute de
        processeur graphique, les plongements ont été calculés sur les
        descriptions écrites et non sur les photographies : deux objets
        semblables mais décrits différemment peuvent échapper au rapprochement.
\end{enumerate}

\conclchap

Ce chapitre a exposé la démarche de l'étude : une recherche interactive de
niveau intégrateur, conduite selon une approche mixte, dont les quatre variables
explicatives ont été rendues mesurables par un ensemble d'indicateurs
observables sur le système lui-même. Il a précisé la population, l'échantillon,
les outils de collecte et les techniques d'analyse, puis la démarche
d'industrialisation retenue. Les limites assumées y ont été énoncées sans
détour, car elles bornent la portée des résultats. Le chapitre suivant présente
le terrain, les données recueillies et les résultats obtenus.
